% `tex.web` §537, §362 and §1335: the `(./file.tex` … `)` a TeX job prints
% around every file it reads.
%
%     print_char("("); incr(open_parens); slow_print(name); update_terminal;
%
% and at a real end of file, §362:
%
%     if name>17 then print_char(")"); decr(open_parens);
%
% and for the files still open when the job stops, §1335:
%
%     while open_parens>0 do
%       begin print(" )"); decr(open_parens); end;
%
% All three are written against the current column, which is why this waited
% for a printer with a column: the `(` joins whatever line is in progress and
% the `)` lands on whatever line is last. That is nearly always §638's page
% echo, so `[0] )` is one line and not two.
%
% The name is the one the file system gave, not the one the document asked
% for — `./sub.tex` for a file beside it, an absolute path for one out of the
% tree — so the resolved path travels back with the contents now.
%
% Every case under tests/ used to be written around the absence of this: a
% report first and last, to keep messages off the lines the echo would have
% been on. `tools/trace_check.sh` still peels the opener, because §252 and
% §299 write `{` with `print_char` too and this engine starts a new line for
% those — that is what is left.
\catcode`\{=1 \catcode`\}=2
\tracingonline=1 \hsize=100pt \vsize=20pt \parindent=0pt
\font\r=cmr10 \r \hbadness=10000 \vbadness=10000
\immediate\openout0=fecho.tex
\immediate\write0{\string\message{inner}}
\immediate\closeout0
\showthe\escapechar

% A file opened and closed: its `)` comes at the real end of the file, and the
% `\message` inside it joins the opener's line.
\input fecho
\message{after}

% ...and a page shipped afterwards, so the echo has something to land on.
A\par
\showthe\escapechar
\end
